\documentclass[10pt,letterpaper]{article}
\usepackage[top=0.85in,left=2.75in,footskip=0.75in]{geometry}
\usepackage{amsmath,amssymb}
\usepackage{changepage}
\usepackage[utf8x]{inputenc}
\usepackage{textcomp,marvosym}
\usepackage{cite}
\usepackage{nameref,hyperref}
\usepackage[right]{lineno}
\usepackage{microtype}
\DisableLigatures[f]{encoding = *, family = * }
\usepackage[table]{xcolor}
\usepackage{array}
\usepackage{listings}
\usepackage{graphicx}
\usepackage{float}
\usepackage{booktabs}
\usepackage{multirow}

% Note: This template is based on PLOS Computational Biology requirements.
% It assumes you have the plos2015.bst file for bibliography if you use \bibliographystyle{plos2015}.
% If not, use \bibliographystyle{unsrt} or similar.

\linenumbers

\title{MedImages.jl: A High-Performance, Differentiable Julia Framework for Medical Image Processing and Metadata Preservation}

\author{Jakub Mitura$^{1*}$, Divyansh Goyal$^{1}$,  Michael C. Kreissl $^{1}$, Joanna Wybranska$^{1}$}

\date{}

\begin{document}

\subsection*{Abstract}

Rationale: Imaging data driven artificial intelligence research in Nuclear Medicine increasingly depends on large-scale multi-modal 3D/4D datasets (Positron Emission Tomography / Single-Photon Emission Computed Tomography (PET/SPECT) with Computed Tomography / Magnetic Resonance Imaging (CT/MRI)) that must be resampled, aligned and quantitatively normalized without losing the spatial and clinical metadata required for physical interpretation (e.g. Standardized Uptake Value (SUV), absorbed dose). Existing pipelines often struggle with (1) biobank-scale volume and input/output (I/O) overhead, (2) execution speed and the ``two-language'' barrier, (3) metadata drift when images are converted to raw arrays (4) and end-to-end differentiability for novel physics-informed machine learning <Jha_Bradshaw_2022_RELAINCE.pdf>.

\textbf{Methods:} We developed \texttt{MedImages.jl}, a native Julia library that binds voxel data to spatial metadata (origin, spacing, direction) and clinical/temporal fields using a typed \texttt{MedImage} structure and a batched multimodal container (\texttt{BatchedMedImage}) for synchronized transforms. Volume-scale performance was evaluated on a 100-subject autoPET PET/CT alignment pipeline (target grid $512^3$) comparing HDF5-based persistence to a Python caching pipeline. The execution speed was benchmarked on Graphics Processing Unit (GPU) for resampling, orientation changes, and fused affine transforms, and compared with SimpleITK (ITK-based C++ reference toolkit <Lowekamp_2013_SimpleITK.pdf>). Differentiability was validated via gradient tests and learned spatial pipelines (inverse-rotation spatial transformer; organ-specific affine registration) and extended to physics-informed dosimetry using a Universal Differential Equation (UDE) for $^{177}$Lu-PSMA trained against Monte Carlo dose ground truths. Metadata integrity was assessed by SUV agreement with SimpleITK and by an end-to-end PET/CT transformation chain with lesion-based readouts.

\textbf{Results:} \texttt{MedImages.jl} achieved a $7.2\times$ turnaround speedup versus caching in the 100-subject benchmark and up to $115\times$ (resampling), $71\times$ (orientation), and $135\times$ (fused affine) GPU speedups versus Central Processing Unit (CPU). SUVbw matched SimpleITK to $10^{-14}$ (double precision), and compounded transforms preserved Mean SUV ($3.4913 \rightarrow 3.4660$; 0.73\%) and lesion volume ($52.55 \rightarrow 51.86$~cm$^3$; 1.32\%), consistent with interpolation effects rather than metadata deterioration. Learned spatial transforms reduced reconstruction error by $>65$\%, and UDE dosimetry refinement converged to Mean Squared Error (MSE) $4.14\times10^{-4}$ versus Monte Carlo ground-truth dose maps <Song_Xue_2024_DoseGAN.pdf>.

\textbf{Conclusions:} MedImages.jl provides a performant, fully differentiable environment for medical image processing. By unifying spatial operations and metadata management within a single language, it facilitates large-scale clinical research and physics-informed deep learning without compromising clinical validity.

\textbf{Keywords:} Medical Imaging, Julia, Scientific Machine Learning, Nuclear Medicine, Differentiable Programming, Spatial Metadata

\section{Introduction}

Nuclear medicine diagnostic data, such as PET and SPECT, provide quantitative measurements of biologic processes that support oncologic staging, assessment of therapy response, and planning of theranostic treatment. Correct processing and interpretation of functional imaging data require that measured uptake is assigned to the same patient-space coordinates as the anatomic reference (CT or MRI). Hybrid imaging therefore combines complementary modalities within a single, physically defined sampling space specified by image origin, grid or voxel spacing, and orientation. 

The clinical landscape of targeted radioligand therapy has been fundamentally altered by the landmark VISION trial, which demonstrated significant improvements in overall survival (OS) and radiographic progression-free survival (rPFS) for patients receiving $^{177}$Lu-PSMA-617 <VISION_Sartor_NEJM_2021.pdf>. This success has led to its regulatory approval and the initiation of a new generation of clinical trials—including PSMAfore <PSMAfore_Lancet_2024.pdf>, SPLASH, ECLIPSE, ProstACT GLOBAL, PSMAddition, and LUNAR—which collectively seek to transition radioligand therapy into earlier pre-chemotherapy and metastatic hormone-sensitive prostate cancer (mHSPC) disease states.

However, as clinical indications for $^{177}$Lu expand, the historical reliance on fixed, empirical dosing regimens (e.g., the standard 7.4~GBq administration per cycle) is increasingly recognized as suboptimal. High variability in inter-patient pharmacokinetics, heterogeneous tumor sink effects, and variable baseline organ functions necessitate a transition toward personalized, dosimetry-guided treatment planning. The therapeutic efficacy of $^{177}$Lu is primarily driven by its $\beta^-$ particle emission, which possesses a relatively short maximum tissue penetration range of approximately 2~mm. This localizes energy deposition but also exerts a localized ``crossfire effect,'' where radiation traverses multiple cell diameters to irradiate adjacent tumor cells that may not directly express the target antigen.

Correct interpretation of functional data requires that measured uptake is assigned to the same patient-space coordinates as the anatomic reference. As multimodal imaging is increasingly reused as input for AI models, managing data volume and maintaining heterogeneous acquisition metadata has become a central practical challenge <Miller_Kuronuma_2022_AI_Chasm.pdf>. A high computational data preparation process typically requires: loading of a large number of examinations from disk, mapping each modality into a shared patient-space coordinate system, resampling onto standardized grids and applying geometric transforms for augmentation or registration. Preserving spatial and quantitative metadata is necessary throughout the multistep preprocessing phase so that voxel values remain physically interpretable. 

Metadata management transcends practical data engineering; it is fundamentally a patient safety and clinical validity issue. The pervasive practice of converting medical image objects from their native, highly regulated formats (such as DICOM or NIfTI) into raw, detached arrays or generic tensors for deep learning consumption routinely decouples the voxel intensities from their strict patient-space definition and from the vital clinical metadata used for quantification <Jha_Bradshaw_2022_RELAINCE.pdf>. This phenomenon, commonly referred to as ``metadata drift,'' can easily yield outputs that appear visually plausible to a neural network but are spatially inconsistent, misaligned, or quantitatively entirely invalid.

In widely used image-processing environments such as SimpleITK <Lowekamp_2013_SimpleITK.pdf> and MONAI <Cardoso_2022_MONAI.pdf>, data handling is often built around caching workflows and chains of separate libraries rather than architectures tailored to biobank-scale diagnostic imaging. As a result, a substantial fraction of total processing time is used not only for core preprocessing operations itself, but also to repeatedly copy very large image volumes between intermediate storage steps and different software components, and to ensure the consistency of spatial information <Dexl_2025_autoPET.pdf>.

\subsection*{Resolving the Nuclear Medicine ``AI Chasm''}
The AI chasm in nuclear medicine refers to the significant gap between the high expectations of artificial intelligence and its actual implementation in clinical practice <Miller_Kuronuma_2022_AI_Chasm.pdf>. This gap is largely driven by a scarcity of high-quality ``AI-ready'' databases, the difficulty of conducting large-scale multi-center evaluations, and the loss of critical quantitative data during preprocessing <Jha_Bradshaw_2022_RELAINCE.pdf>.

MedImages.jl helps resolve this chasm through several key innovations:
\begin{itemize}
    \item \textbf{Preventing ``Metadata Drift'' for Quantitative Fidelity}: A major barrier to clinical AI is that converting medical image objects into raw arrays for deep learning often detaches voxel intensities from their patient-space definitions and clinical metadata. MedImages.jl solves this by binding voxel data to spatial metadata (origin, spacing, direction) and clinical fields (injection time, patient weight) using a typed \texttt{MedImage} structure. This ensures that critical physical interpretations—such as Standardized Uptake Value (SUV) and absorbed radiation dose—are preserved exactly at machine precision (matching reference standard tools to $10^{-14}$).
    \item \textbf{Biobank-Scale Scalability}: Bridging the AI chasm requires massive, multi-center datasets <Miller_Kuronuma_2022_AI_Chasm.pdf>. MedImages.jl utilizes a fused affine preprocessing strategy combined with HDF5-based storage and GPU acceleration to process data at extreme speeds. By collapsing spatial transformations (normalization, resampling, centering) into a single computational step, it achieves up to a $135\times$ GPU speedup and reduces per-subject turnaround time to less than 200 milliseconds, solving the I/O and processing bottlenecks of biobank-scale workflows.
    \item \textbf{Synchronized Multimodal Alignment}: AI in nuclear medicine heavily depends on hybrid datasets (e.g., PET/CT, SPECT/CT). MedImages.jl uses a \texttt{BatchedMedImage} container that applies identical, synchronized geometric transformations across multiple modalities simultaneously without manual bookkeeping. This guarantees strict voxel-wise spatial correspondence across input channels, which is a fundamental prerequisite for training reliable multimodal neural networks.
\end{itemize}

\subsection*{Enabling Biology- and Physics-Informed Neural Networks (SciML)}
Scientific Machine Learning (SciML)—which intertwines known physical and biological laws with deep learning—is a major frontier for precision care, such as creating ``theranostic digital twins'' <Miller_Kuronuma_2022_AI_Chasm.pdf>. MedImages.jl accelerates this by making geometric and spatial operations natively differentiable.
\begin{itemize}
    \item \textbf{End-to-End Differentiable Pipelines}: Because MedImages.jl is built natively in Julia, it integrates seamlessly with automatic differentiation frameworks like \texttt{Zygote.jl} and \texttt{Enzyme.jl}. This means that geometric image operations (resampling, rotation, cropping) can be incorporated directly into end-to-end model optimization without breaking the gradient flow or requiring non-differentiable C++ backends.
    \item \textbf{Triple-Branch Universal Differential Equations (UDEs) for Dosimetry}: MedImages.jl facilitates high-fidelity, physics-informed dosimetry by partitioning mechanistic physical laws from learned stochastic scattering. By unifying hardware acceleration, absolute metadata integrity, and differentiable physics within a single language stack, MedImages.jl bridges the divide between the physics of image acquisition and the logic of AI inference.
\end{itemize}


In this work, we evaluate MedImages.jl across nuclear medicine–relevant benchmarks and experiments, both in isolation and in integration with complementary Julia libraries, with emphasis on (i) scalability for large PET/CT cohorts with high preprocessing demands, (ii) execution speed without the maintenance costs of a two-language architecture, (iii)  metadata preservation to maintain the physical interpretability required for quantitative nuclear medicine, and (iv) differentiability of geometric operations for learned spatial transforms.


\section{Methods}

\begin{figure}[H]
\centering
\includegraphics[width=1.0\linewidth]{../article_sources/new_images/graphical_abstract.png}
\caption{End-to-End Multi-Modal SciML Dosimetry Pipeline. Phase 1 shows raw clinical modalities in coronal view with typical acquisition misalignments (rotation and scaling discrepancies). Phase 2 demonstrates the result of unified preprocessing and batch alignment on a shared grid with soft-tissue windowed CT. Phase 3 illustrates the integration of these data into the Universal Differential Equation (UDE) framework. Phase 4 displays the final dosimetry inference and error residuals compared to Monte Carlo ground truth.}
\label{fig:flowchart}
\end{figure}

MedImages.jl was developed to address several practical requirements that are highly relevant in medical imaging, including preservation of image geometry and acquisition metadata, reliable handling of clinical image formats, efficient processing of large multimodal datasets, and support for optimization-based learning workflows. 
A central design feature is the \verb|MedImage| data structure, which stores voxel values together with essential spatial information such as image position, voxel spacing, orientation, patient identity, and acquisition time. This avoids a common problem in AI-oriented workflows, where images are converted into generic tensors and thereby lose information needed to interpret them in physical space \cite{}. 
For robust reading and writing of clinical formats such as DICOM and NIfTI, MedImages.jl uses ITK through a wrapper interface, whereas spatial operations such as rotation, translation, and resampling are implemented directly in Julia language. This design combines the reliability of a mature medical imaging toolkit for file handling with the flexibility to use spatial operations in modern learning and optimization pipelines. 
For larger studies, the \verb|BatchedMedImage| structure allows multiple images to be processed together while preserving the spatial properties of each individual volume. In combination with HDF5-based storage and GPU acceleration, this architecture reduces data-loading overhead, improves scalability, and supports workflows in which image transformations can be incorporated directly into model optimization.  

Moreover, compatibility with the broader Julia SciML ecosystem allows MedImages.jl to interoperate with other scientific computing libraries and supports differentiable geometric operations for learned spatial transforms.  

To demonstrate that MedImages.jl addresses the central computational and metadata-related challenges of large-scale medical image analysis, we conducted three proof-of-concept experimental series. An overview of the end-to-end multi-modal pipeline is illustrated in Figure \ref{fig:flowchart}, and the core framework functionalities are summarized in Table \ref{tab:medimages_functionalities}.

\begin{table}[H]
\centering
\caption{Core functionalities and data structures of MedImages.jl for 3D medical image processing.}
\label{tab:medimages_functionalities}
\renewcommand{\arraystretch}{1.2}
\begin{tabular}{|p{0.24\textwidth}|p{0.68\textwidth}|}
\hline
\textbf{Functionality} & \textbf{Description} \\
\hline
\texttt{MedImage} data structure & Unified medical image representation that stores voxel data together with clinical and vectorized spatial metadata. \\
\hline
\texttt{BatchedMedImage} structure & Batched representation for joint processing of multiple images while preserving the spatial properties of each individual volume. \\
\hline
Imaging-format-compatible I/O & Loading of standard medical image formats, including NIfTI and DICOM, via \texttt{ITKIOWrapper}. \\
\hline
HDF5 support & Fast loading and saving in big-data compatible format, facilitating efficient intermediate storage in repeated processing and deep-learning workflows. \\
\hline
GPU acceleration & Transparent GPU support for spatial transformations through \texttt{KernelAbstractions.jl}, enabling accelerated parallel execution of preprocessing and geometric operations. \\
\hline
Resampling & Native Julia functions for resampling, enabling voxelwise correspondence across modalities. \\
\hline
Spatial transformations & Native Julia functions for manipulation of spatial properties of 3D image volumes, including \texttt{rotate\_mi}, \texttt{translate\_mi}, \texttt{scale\_mi}, \texttt{resample\_to\_spacing}, \texttt{crop\_mi}, and \texttt{pad\_mi}. \\
\hline
Orientation management & Tools for handling, preserving, and converting image orientations across neuroimaging and medical imaging coordinate systems. \\
\hline
Automatic differentiation & Differentiable implementation of spatial transformations with well-defined gradients through \texttt{Zygote.jl} and \texttt{Enzyme.jl}. \\
\hline
\end{tabular}
\end{table}




\textbf{Experimental Series 1: Evaluation of Scalability and Execution Speed}

For large-scale imaging studies, we evaluated and compared MedImages.jl against contemporary frameworks with respect to two key requirements: scalability to massive multimodal cohorts and raw runtime of core spatial preprocessing operations in repeated deep-learning workflows.

To assess scalability and computational overhead, we designed a 100-case preprocessing benchmark based on the autoPET dataset <Dexl_2025_autoPET.pdf>. This setup reflects a realistic biobank-scale scenario where serial PET/CT studies must be repeatedly loaded and geometrically prepared during iterative model training.

The main preprocessing task was multimodal spatial alignment of heterogeneous PET and CT volumes on a common isotropic reference grid of $512^3$ voxels. This harmonization is essential for multimodal neural networks, which require strict voxelwise spatial correspondence across input channels.

This benchmark was divided into two computational phases per subject: first, loading the preprocessed case from cache, and second, applying the remaining transformations required to prepare the tensor data for training.

Within the MedImages.jl pipeline, HDF5-based storage was combined with a fused affine preprocessing strategy to reduce computational overhead. Rather than performing orientation normalization, voxel spacing adjustment, and spatial centering as separate steps, these operations were merged into a single affine transformation Pass.

MedImages.jl was compared with MONAI <Cardoso_2022_MONAI.pdf> using its PersistentDataset mechanism, thus assessing not only the framework-level performance, but also the effect of different data persistence strategies. Performance was evaluated under repeated training conditions, beginning from epoch 2 onward.

To assess raw computational speed independently of the caching benchmark, isolated spatial operations were measured across multiple scales. These tests included resampling, orientation changes, and affine transformations. Benchmarks were performed on single volume of size $256 \times 256 \times 128$ using both CPU and GPU backends. SimpleITK <Lowekamp_2013_SimpleITK.pdf> was evaluated on CPU, reflecting its typical use, whereas MedImages.jl was evaluated on both CPU and GPU to quantify the hardware acceleration effect.

Both the 100-subject caching pipeline benchmark and the isolated raw computational speed benchmarks were executed on strictly identical hardware: an Intel Core Ultra 9 285K processor and an NVIDIA RTX 3090 GPU.

 . 

\textbf{Experimental Series 2: Verification of Metadata Integrity and Quantitative Clinical Fidelity  }

To evaluate the quantitative accuracy of MedImages.jl, Standardized Uptake Value normalized by body weight (SUVbw) was compared against a reference implementation in SimpleITK <Lowekamp_2013_SimpleITK.pdf>. SUVbw was computed using same PET image data from autopet dataset and corresponding clinical and temporal metadata in both frameworks. The comparison was performed in double-precision floating point arithmetic. For the fused affine interpolation kernel, the numerical results obtained on the CPU and GPU were compared directly, and numerical agreement was strictly verified within a tolerance of $10^{-4}$ to guarantee floating-point consistency across backends.


An end-to-end consistency experiment was subsequently conducted on a subset of autoPET dataset to assess the preservation of quantitative fidelity under a sequence of spatial transformations. The processing pipeline consisted of (1) segmentation of lesion masks on the original image data, (2) conversion of PET intensities to SUV, (3) resampling of PET volumes and corresponding lesion masks to the native CT coordinate space, (4) resampling to an isotropic voxel spacing of 2.0 mm, (5) application of a $45^\circ$ rotation around the z-axis, and (6) translation by 10 voxels along the x-axis. Nearest-neighbor interpolation was applied to segmentation masks, while linear interpolation was used for continuous PET and CT intensity data. Quantitative evaluation was performed by measuring mean SUV and total lesion volume within the segmentation mask before and after the transformation sequence (see Figure \ref{fig:transformations}).

\begin{figure}[H]
\centering
\includegraphics[width=1.0\linewidth]{../article_sources/new_images/transformations.png}
\caption{End-to-end consistency experiment illustrating the sequence of spatial transformations applied to the autoPET dataset. The pipeline demonstrates the preservation of quantitative fidelity during sequential spatial manipulation steps including conversion to SUV, resampling to native CT space, isotropic resampling, rotation, and translation.}
\label{fig:transformations}
\end{figure}

To evaluate the capability of MedImages.jl for handling complex theranostic imaging data, a multi-modal processing experiment was conducted using a dataset consisting of four imaging modalities from a single patient who underwent a radioligand therapy. The data set included: (1) [$^{177}$Lu]Lu-PSMA SPECT with attenuation correction (AC) (2) [$^{177}$Lu]Lu-PSMA SPECT without attenuation correction (NAC), (3) a voxel-wise absorbed dose map, and (4) a CT volume serving as anatomical reference. Each modality was assumed to originate from independent acquisition and reconstruction pipelines and therefore retained its own voxel grid, data type, and spatial metadata, including voxel spacing, image origin, and orientation. All four modalities were loaded into a single BatchedMedImage data structure. During this step, the original voxel data and modality-specific spatial metadata were preserved without prior resampling or harmonization. Subsequently, a geometric transformation was applied to the entire batch. A spatial rotation of $45^{\circ}$ around the Z-axis (axial plane) was executed simultaneously across all four modalities using a single batched transformation call, followed by a linear translation of 10 voxels along the X-axis to verify spatial origin tracking.


The outcome of the transformation was evaluated qualitatively by visual inspection of corresponding axial slices before and after transformation. Particular attention was given to the preservation of spatial alignment between functional tracer distribution and anatomical structures. The consistency of spatial relationships across modalities after transformation served as the primary validation criterion. The exact methodology parameters for this experiment are detailed in Table \ref{tab:method_info_batched_Lu177} in the Appendix. 

 All experiments described are open-source and reproducible via the codebase provided in this repository.

Experimental Series 3: Verification of Differentiability and Learned Spatial Transformations

To determine whether spatial image operations remained mathematically trainable and could therefore be used reliably in end-to-end learning pipelines, gradient propagation was numerically evaluated across representative geometric transformations implemented in MedImages.jl. Reverse-mode gradients were computed with Zygote.jl for translation, padding, cropping, and fused affine interpolation. For translation and padding, correctness was assessed by verifying that the gradient remained constant throughout the transformed image domain. For cropping, the gradient pattern was examined to confirm preservation within the retained image region and suppression outside the cropped area. For the fused affine interpolation kernel, the gradients obtained on the CPU and GPU were compared directly, and numerical agreement was verified within a tolerance of $10^{-4}$. 

To further assess end-to-end differentiability in a trainable setting, a learned inverse rotation experiment was conducted using synthetic single-case batched image configurations. A synthetic 3D structural line image was generated on a strictly confined $16 \times 16 \times 16$ voxel grid to ensure strong spatial gradient signals. A total of 120 batched sample configurations (100 for training, 20 for testing) were generated by applying random three-axis sequential rotations within an angular range of $\pm 30^{\circ}$. A lightweight 3D Convolutional Neural Network (CNN) combined with a Multi-Layer Perceptron (MLP) head was then trained using gradient descent for 20 epochs to predict the inverse Euler rotation parameters directly from the spatial grids.

\section{Results }

All results in scalability, computational performance, metadata integrity, quantitative fidelity, and differentiability are summarized in Table~\ref{tab:all_results}.

\subsection*{Evaluation of Scalability and Processing Speed}
In iterative training scenarios, MedImages.jl demonstrated substantially reduced processing times compared to established frameworks, decreasing total turnaround per subject from several hundred milliseconds (MONAI <Cardoso_2022_MONAI.pdf>) to well below one-tenth of a second. The greatest improvements were observed during the transformation stage, where the fused affine pipeline reduces repeated interpolation and data movement by performing multiple spatial operations in a single computational step. 

GPU acceleration further amplified performance gains across core spatial operations. For medium-sized volumes ($256 \times 256 \times 128$), resampling, orientation changes, and affine transformations achieved speedups exceeding one to two orders of magnitude compared with CPU execution, with additional improvements over standard CPU-based libraries. Notably, while frameworks like MONAI also utilize GPU acceleration for transformations, their typical implementation relies on a sequence of individual operations (e.g., orientation normalization followed by resampling). In our benchmarks, this sequential approach resulted in a computation time of approximately 500 ms per subject on the GPU. In contrast, MedImages.jl's fused affine kernel mathematically collapses the entire transformation stack into a single Pass, achieving a $10\times$ computational speedup on the same hardware. In large-scale settings ($512^3$ target grid, $n=100$), per-subject processing times remained below 200~ms, confirming suitability for high-throughput and biobank-scale workflows.

\subsection*{Evaluation of Metadata Integrity and Quantitative Fidelity}
Quantitative validation demonstrated numerical agreement with the reference standard at machine precision for SUV calculations, with differences on the order of $10^{-14}$. Across compound spatial transformations, mean SUV changed only minimally (from 3.49 to 3.47), and lesion volume remained stable (from approximately 52.6~cm$^3$ to 51.9~cm$^3$), corresponding to relative deviations below 1--2\%. These differences are consistent with expected interpolation effects and do not indicate degradation of clinical metadata.

\subsection*{Evaluation of Differentiability and Learned Transformations}
Gradient propagation across spatial operations followed expected analytical behavior, with identity-preserving transformations maintaining uniform gradients and spatially selective operations producing region-dependent responses. Numerical agreement between CPU and GPU gradient computations was maintained within a tolerance of $10^{-4}$.

In learned transformation experiments, optimization of geometric parameters resulted in a reduction in reconstruction error exceeding 65\% within 20 training iterations, demonstrating effective end-to-end training. Affine registration tasks showed stable convergence using GPU-based gradient optimization, confirming that complex spatial transformations can be integrated directly into differentiable learning pipelines without external implementations. Figure \ref{fig:rotation} visualizes the learned inverse rotation trajectory.

\begin{figure}[H]
\centering
\includegraphics[width=0.8\linewidth]{../article_sources/new_images/differentiability.png}
\caption{Learned inverse rotation trajectory across training epochs, demonstrating end-to-end spatial differentiability. The model progressively aligns the uncorrected volume (red) towards the gold standard (blue).}
\label{fig:rotation}
\end{figure}

\begin{table}[H]
\centering
\caption{Summary of results across all experimental series.}
\label{tab:all_results}
\renewcommand{\arraystretch}{1.2}
\begin{tabular}{|l|l|c|c|}
\hline
\textbf{Experiment Series} & \textbf{Metric / Task} & \textbf{Result} & \textbf{Reference / Comparison} \\
\hline

\multicolumn{4}{|c|}{\textbf{Series 1: Scalability and Execution Speed}} \\
\hline
Data loading & Load from cache & \textbf{~40 ms} & MONAI: ~150 ms \\
Computation & Transform (compute) & \textbf{~50 ms} & MONAI: ~500 ms\\
Overall & Total turnaround & \textbf{~90 ms} & MONAI: ~650 ms\\
GPU performance & Resampling & \textbf{115$\times$ faster (GPU vs CPU)} & 17$\times$ vs SimpleITK \\
GPU performance & Orientation changes & \textbf{71$\times$ faster (GPU vs CPU)} & 31$\times$ vs SimpleITK \\
GPU performance & Fused affine transform & \textbf{135$\times$ faster (GPU vs CPU)} & 8.1$\times$ vs SimpleITK \\
Large-scale pipeline & $512^3$ volumes (100 cases) & \textbf{< 200 ms per subject} & High-throughput feasible \\
\hline

\multicolumn{4}{|c|}{\textbf{Series 2: Metadata Integrity and Quantitative Fidelity}} \\
\hline
SUV validation & SUVbw agreement & \textbf{$\approx 10^{-14}$} & Matches SimpleITK (machine precision) \\
Transformation consistency & Mean SUV deviation & \textbf{0.73\%} & Before vs after pipeline \\
Transformation consistency & Lesion volume deviation & \textbf{1.32\%} & Before vs after pipeline \\
\hline

\multicolumn{4}{|c|}{\textbf{Series 3: Differentiability and Learned Transformations}} \\
\hline
Gradient validation & Identity transforms & \textbf{Gradient = 1} & Constant gradient propagation \\
Gradient validation & Cropping & \textbf{1 (inside), 0 (outside)} & Expected spatially selective gradient pattern\\
Gradient validation & CPU vs GPU consistency & \textbf{$< 10^{-4}$ difference} & Numerical consistency between hardware backends\\
Learning experiment & Inverse rotation (CNN) & \textbf{> 65\% MSE reduction} & End-to-end optimization\\
Registration & Affine optimization & \textbf{Converged} & GPU-based gradient optimization \\
Improved UDE dosimetry & Final voxel-wise dose error & \textbf{MSE = $4.14 \times 10^{-4}$} & Monte Carlo ground-truth dose maps\\
\hline

\end{tabular}
\end{table}


\section{High-Fidelity Dosimetry via Triple-Branch Universal Differential Equations}

\subsection{Model Architecture: The Triple-Branch UDE}
The proposed model evolves the standard dosimetry framework into a Universal Differential Equation (UDE) that partitions mechanistic physical laws from learned stochastic scattering. The core innovation is the expansion to a \textbf{Triple-Branch Input} architecture:
\begin{enumerate}
    \item \textbf{Activity Branch ($A_{total}$)}: Tracks the biological time-varying radiopharmaceutical concentration derived from the compartmental ODE states (free and bound activity).
    \item \textbf{Density Branch ($\rho$)}: Incorporates voxel-specific mass-density derived from CT data.
    \item \textbf{Gradient Branch ($|\nabla \rho|$)}: Acts as a geometric prior to explicitly signal tissue-bone and tissue-air boundaries, which allows the neural network to learn asymmetric Compton scattering.
\end{enumerate}

\subsection{Numerical Stability: Balanced Unit Training}
To ensure convergence on GPU architectures, we implemented \textbf{Balanced Unit Training}. SPECT activity variables are internally re-scaled by $10^{-6}$ to bring the IVP (Initial Value Problem) into an $O(1)$ range. The physical dose constant ($C_{conv}$) is balanced accordingly, ensuring the integrated output maintains its physical identity in Grays (Gy).

\subsection{Physical Formulation}
The system is defined by the following coupled system of equations:
\begin{equation}
    \frac{d}{dt} \text{Dose}(t) = \text{softplus} \left( \frac{A_{total}(t) \cdot C_{conv}}{\rho \cdot V} + \mathcal{N}_\theta(A_{std}, \rho_{std}, |\nabla \rho|_{std}) \right)
\end{equation}
Where $\mathcal{N}_\theta$ is a 3D Convolutional Neural Network learning the corrective residual to the analytical dose rate. 

The formulation integrates three distinct computational logic blocks:
\begin{enumerate}
    \item \textbf{Integrating Classical Physics with Deep Learning}: The software formulates energy deposition analytically (calculating dose rate via $\frac{A_{total}(t) \cdot C_{conv}}{\rho \cdot V}$) while simultaneously using a 3D Convolutional Neural Network ($\mathcal{N}_\theta$) as a learned corrective residual for complex stochastic effects.
    \item \textbf{Neural Network Residual}: The term $\mathcal{N}_\theta(A_{std}, \rho_{std}, |\nabla \rho|_{std})$ acts as a learned correction for stochastic effects such as asymmetric Compton scattering. By utilizing a Triple-Branch input (standardized activity, density, and the density gradient $|\nabla \rho|$), the model explicitly identifies tissue boundaries where classical analytical kernels typically fail.
    \item \textbf{Differentiable Positivity}: The entire rate is wrapped in a \textit{softplus} activation function. This ensures that the predictions remain physically non-negative while providing a smooth, continuously differentiable surface required for stable backpropagation through the ODE solver during end-to-end SciML training.
\end{enumerate}
By unifying hardware acceleration, absolute metadata integrity, and differentiable physics within a single language stack, MedImages.jl bridges the divide between the physics of image acquisition and the logic of AI inference.


\subsection{Results and Comparative Analysis}
The Improved UDE was evaluated on a full-body validation cohort using a sliding-window reconstruction ($64^3$ patches, stride 32). The performance of the Triple-Branch UDE was benchmarked against traditional static physics models, state-of-the-art pure deep learning architectures (like DblurDoseNet <Li_2022_DblurDoseNet.pdf>), and the gold-standard Monte Carlo (MC) radiation transport simulations <Mansouri_2024_Transformer_Dosimetry.pdf>. Quantitative performance metrics are summarized in Table \ref{tab:dosimetry_performance}, and a qualitative comparison of the corresponding dose maps is presented in Figure \ref{fig:comparison_dosemaps}.

\begin{table}[h]
\centering
\caption{Full-Body Performance Comparison (Pearson Correlation and Mean Absolute Error)}
\label{tab:dosimetry_performance}
\begin{tabular}{lll}
\hline
\textbf{Model} & \textbf{Mean Pearson ($r$)} & \textbf{MAE (Gy)} \\ \hline
Analytical Baseline (Static) & 0.8249 & 700.43 \\
DblurDoseNet (Pure DL) & 0.5566 & 1102.34 \\
Spect0Net & 0.6124 & 955.12 \\
SemiDose & 0.6401 & 912.88 \\
Original UDE (2-Branch) & 0.9031 & 405.77 \\
\textbf{UDE IMPROVED (Triple-Branch)} & \textbf{0.9221} & \textbf{359.49} \\ \hline
\end{tabular}
\end{table}

\begin{figure}[H]
\centering
\includegraphics[width=1.0\linewidth]{../article_sources/new_images/comparison_dosemaps.png}
\caption{Multi-model comparative dosimetry. Qualitative evaluation of the proposed Triple-Branch UDE framework against purely analytical baseline and deep-learning (CNN) models. Error residuals are computed relative to the gold-standard Monte Carlo simulated absorbed dose maps, demonstrating superior performance and reduced hallucination artifacts.}
\label{fig:comparison_dosemaps}
\end{figure}

\subsection{Comparative Efficacy and Clinical Viability}
To objectively evaluate the UDE model, it must be contrasted with the prohibitive computational latency of the gold-standard Monte Carlo simulations. While MC particle tracking represents the absolute pinnacle of physical accuracy, generating a high-fidelity, low-noise dose map for a single single-bed SPECT study requires approximately \textbf{48~hours (2 days)} of continuous processing on standard clinical CPU infrastructure <Kim_2024_CoRX_NET.pdf>. This creates an insurmountable bottleneck for real-time adaptive therapy <Mansouri_2024_Transformer_Dosimetry.pdf>.

In contrast, while pure convolutional networks like DblurDoseNet can perform inference in approximately 3~seconds <Li_2022_DblurDoseNet.pdf>, they operate as ``black box'' models prone to catastrophic hallucination in out-of-distribution (OOD) scenarios, such as when crystal thickness or reconstruction algorithms shift across centers. The UDE framework establishes a middle ground: achieving near real-time inference (measured in seconds to tens of seconds) while rigorously anchoring predictions in exact physical laws. 

The UDE solves the inherent ``black box'' crisis of modern artificial intelligence in healthcare. By explicitly programming the known mechanisms of radiation decay and fluid compartment transport into the system, the model structurally prevents physiological violations. It relegates the neural network to its optimal role: approximating the complex, non-linear patient heterogeneities—such as unique tumor micro-environment clearance rates—that simply cannot be described by classical analytic functions <Akhavanallaf_2021_PINN.pdf>. 

Because the UDE is irrevocably anchored by invariant physical laws, its parameter space is vastly reduced and highly constrained compared to unconstrained models like DblurDoseNet. This physics-guided constraint makes UDEs exceptionally robust to dataset shifts, a monumental advantage in heterogeneous multi-center clinical trials <Song_Xue_2024_DoseGAN.pdf>. Furthermore, the UDE's ability to extrapolate accurate multi-day time-activity curves (TACs) from limited boundary conditions makes it the only theoretically sound choice for Single-Time-Point (STP) dosimetry, which is rapidly becoming the desired clinical standard to minimize patient burden. 

Finally, by partitioning mechanistic laws from learned biological clearance rates ($\mu_{bio}$), the UDE model provides a level of interpretability and regulatory trust absent in pure DL architectures <Akhavanallaf_2021_PINN.pdf>. Final outputs can be decomposed into their constituent physical and learned components, allowing medical physicists to interrogate the biological parameters for physiological plausibility (e.g., ensuring non-negative clearance rates) before clinical decision-making. The proposed paradigm thus stands as the apex model for personalized radiopharmaceutical dosimetry in the modern theranostic era.


\subsection{Computational Performance: SciML vs. Pure Deep Learning}
While the Triple-Branch UDE achieves superior accuracy through physics-informed constraints, integrating a system of ordinary differential equations (ODEs) inherently increases computational complexity compared to purely feed-forward deep learning models. 

To properly benchmark the UDE approach and address the multi-language landscape, we implemented mathematically identical Triple-Branch UDEs (same 3D CNN architectures and ODE formulations) across three frameworks: Native Julia (SciML + Lux.jl + Zygote), PyTorch (\texttt{torchdiffeq}), and JAX (\texttt{diffrax} via Equinox). We additionally included a hybrid Python approach (\texttt{scipy.integrate.solve\_ivp} interacting with a PyTorch CNN). To ensure isolation of purely computational capabilities and to allow for fully open-source reproducibility without relying on protected patient data, all framework-level execution and memory benchmarks were performed using synthetic random tensors simulating a $32^3$ voxel patch. All tests were run on an NVIDIA H100 PCIe GPU, explicitly excluding Just-In-Time (JIT) compilation overhead to ensure absolute equivalency.

\textbf{Inference Speed (300-hour integration):}
\begin{itemize}
    \item \textbf{JAX (\texttt{diffrax})}: \textbf{9.4 ms} (Highest performance due to end-to-end XLA compilation of the ODE solver and neural network).
    \item \textbf{PyTorch (\texttt{torchdiffeq})}: \textbf{24.0 ms}.
    \item \textbf{Julia SciML (\texttt{Euler})}: \textbf{48.4 ms} (Highly competitive native performance without requiring strict XLA constraints).
    \item \textbf{Hybrid Python (\texttt{scipy.integrate})}: \textbf{675.2 ms} (Severely bottlenecked by CPU-GPU data transfer overheads).
\end{itemize}

\textbf{Training Performance (Forward + Backward Pass) \& Memory Analysis:}
Inspired by the architectural evaluations in \texttt{SciMLBenchmarks.jl}, we conducted a deeper analysis into the memory scaling and adjoint sensitivity methods utilized by the different frameworks during backpropagation through the ODE solver:
\begin{itemize}
    \item \textbf{JAX}: \textbf{52.7 ms} per step. Peak memory: \textbf{197.4 MB}. (JAX achieves extraordinary efficiency by compiling the entire reverse-mode ODE solve into a single optimized XLA graph).
    \item \textbf{PyTorch}: \textbf{62.9 ms} per step. Peak memory: \textbf{1.37 GB}. (High memory consumption is observed due to the standard practice of retaining the large 3D CNN computational graph across all ODE solver steps during backpropagation).
    \item \textbf{Julia SciML (BacksolveAdjoint)}: \textbf{478.8 ms} per step. Peak memory: \textbf{$\approx$ 2.45 GB}. (By applying insights from \texttt{SciMLBenchmarks.jl} and using a \texttt{BacksolveAdjoint}, the solver re-integrates the ODE backward in time instead of storing dense checkpoints. Slower training times currently reflect the overhead of \texttt{Zygote.jl} tracking complex neural network VJPs within the ODE solver loop, compared to the mature C++ backends of PyTorch/JAX).
    \item \textbf{Julia SciML (EnzymeVJP Trial)}: Following best practices from \texttt{SciMLBenchmarks.jl} for neural ODEs, we attempted to compile the reverse-mode ODE adjoint directly using \texttt{Enzyme.jl}. While Enzyme provides unmatched performance for scalar and small-scale array codes, attempting to JIT-compile the reverse-mode AD graph for a 3D neural ODE on the GPU resulted in timeouts and hard CUDA memory crashes after over an hour of compilation. This empirically highlights that while theoretically superior, compiling highly-branched, large-scale Neural PDEs with Enzyme on GPUs currently exceeds practical hardware compiler limits, remaining an active area of development.
\end{itemize}

Although the native Julia UDE approach is slower than JAX/XLA in this specific fixed-step, GPU-bound training setup, its inference time (48.4 ms per patch) remains highly competitive and easily scalable to full-body reconstruction (under 3 minutes per patient). It is important to contextualize this performance within the broader landscape of differential equation solvers. 

As extensively documented by the \texttt{SciMLBenchmarks.jl} suite (specifically the \textit{Multi-Language Wrapper Benchmarks} comparing Julia, SciPy, MATLAB, and deSolve), native Julia ODE solvers systematically outperform multi-language wrapper packages by orders of magnitude on standard stiff and non-stiff problems. In general cases, SciML's dramatic superiority stems from the elimination of cross-language interpreter overhead (e.g., Python repeatedly calling C++ solver backends), the ability to aggressively JIT-compile the entire solver loop, and the use of highly optimized stack-allocated arrays for small-to-medium physical systems.

Our specific neural dosimetry case diverges from these general benchmarks because the computational load is overwhelmingly dominated by massive, static 3D convolutions on the GPU rather than the ODE solver's internal stepping logic. This architectural pattern---a fixed-step solver wrapping a static, massive tensor graph---is uniquely suited to JAX's end-to-end XLA compiler, which fuses the entire operation into a monolithic GPU kernel, yielding the observed 9.4 ms inference speed. 

However, based on the fundamental limits exposed by the SciML multi-language benchmarks, Julia SciML's performance superiority is projected to re-emerge rapidly as the dosimetry model approaches physiological reality. When the biological system becomes heavily stiff (e.g., modeling rapid receptor saturation, non-linear pharmacokinetics, or instantaneous drug injection pulses requiring discrete callbacks), simple fixed-step solvers like Euler fail numerically, and advanced adaptive-step solvers are required. Crucially, the Julia implementation remains the only one capable of natively handling the stiff versions of these equations (e.g., seamlessly switching to advanced solvers like \texttt{KenCarp4} or \texttt{Rodas5}) without requiring bespoke C++ wrappers. 

Implementing and differentiating through complex, dynamically branching stiff solvers in JAX or PyTorch is fundamentally prohibitive without writing custom CUDA kernels. Thus, the MedImages.jl and SciML framework trades raw feed-forward speed on static GPU graphs for unparalleled physical accuracy and solver flexibility, which is critical for maintaining the physical validity of complex biological dosimetry models where concentration gradients shift abruptly.

\subsection{Limitations and Simplifications}
The primary objective of this experiment was to serve as a proof of concept, demonstrating the validity and feasibility of the UDE-based approach in Julia. While the UDE model definitively achieves the highest accuracy among the tested baselines, several physical abstractions were employed to simplify the initial implementation. To maintain scientific rigor and satisfy future clinical peer review, we acknowledge the following physical realities that the current framework abstracts:
1. \textbf{Static Anatomy}: The model assumes a rigid day-0 CT geometry, neglecting 4D respiratory motion.
2. \textbf{Isotropic PK}: Diffusion between neighboring voxels is currently ignored in the interstitial compartment.
3. \textbf{Receptor Dynamics}: Receptor capacity ($B_{MAX}$) is treated as a constant, neglecting radiation-induced receptor degradation (``stunning'').
4. \textbf{Transport Lumping}: Local $\beta^-$ and penetrating $\gamma$ transport are lumped into a single neural correction term rather than a dual-kernel separation.

Despite these necessary simplifications, the experiment successfully proves the validity of the approach: intertwining known physics with deep learning via Julia's SciML stack yields substantially superior dosimetry results compared to pure deep learning models and classical analytical convolutions.

\section{Discussion}

As biomedical imaging increasingly embraces complex multimodal, multidimensional, and longitudinal studies, alongside scientific machine learning (SciML), the interoperability and performance of data management frameworks become critical. Our results indicate that MedImages.jl’s metadata handling paradigm, combined with seamless integration into the Julia scientific ecosystem, provides a highly performant and differentiable framework that systematically addresses the challenges of large-scale clinical AI.

\subsection*{Addressing the Speed and Volume Bottlenecks}
Accurate metadata alignment is a fundamental requirement for contemporary multimodal image analysis. For example, training convolutional neural networks (CNNs) on combined PET/CT data requires multi-channel tensors in which voxel (i,j,k) in the PET dataset corresponds precisely to the same anatomical location as corresponding voxel in the CT dataset. Owing to the large size of 3D voxel-based imaging data, the necessary continuous spatial resampling may represent the principal computational bottleneck in deep learning workflows. 

Integration of metadata within the \texttt{MedImage} structure directly translates into performance gains observed in our benchmarks. The 7.2$\times$ turnaround advantage observed in high-frequency training loops is primarily attributed to MedImages.jl's integration with HDF5—which bypasses the serialization costs inherent in many Python caching mechanisms—and its single-pass Fused Affine kernel. Notably, while established frameworks like MONAI <Cardoso_2022_MONAI.pdf> also offer GPU-accelerated transformations, they typically apply operations as a sequence of separate kernels (e.g., orientation followed by resampling). By mathematically collapsing the entire spatial transformation stack into a single Pass, MedImages.jl eliminates the memory traffic and synchronization overhead associated with sequential operation queues, providing a highly efficient solution for biobank-scale deployments.

\subsection*{Beyond the Python-C++ Divide: Differentiable Imaging Physics}
The Julia ecosystem is uniquely positioned for SciML because its automatic differentiation frameworks (\texttt{Zygote.jl}, \texttt{Enzyme.jl}) can differentiate through arbitrary native code, including custom spatial transformations. Nuclear medicine imaging is inherently quantitative, but its key readouts (e.g., SUV, absorbed dose) are highly sensitive to geometric preprocessing such as resampling and registration. In practice, these operations must often be integrated into AI pipelines where transform parameters are optimized jointly with network weights. However, many commonly used imaging backends execute spatial resampling outside the automatic differentiation graph, preventing end-to-end training through geometric steps. MedImages.jl enables fully differentiable volumetric spatial operations, allowing geometric image operations (resampling, rotation, cropping) to be incorporated directly into model optimization without breaking the gradient flow.

\subsection*{Comparison with Modern Frameworks}
Modern Python frameworks like MONAI <Cardoso_2022_MONAI.pdf> have introduced sophisticated mechanisms to couple metadata with image arrays. MONAI utilizes a \texttt{MetaTensor} structure, appending a metadata dictionary to the PyTorch Tensor. A similar design pattern appears in TorchIO, which prioritizes convenient multimodal coordination but still builds its core processing around external imaging backends. This dependency chain maintains a layer of indirection that can lead to memory management issues in large-scale training.

In contrast, MedImages.jl leverages Julia's type system to bind metadata directly within the \texttt{MedImage} struct. By removing the runtime overhead associated with dictionary-based metadata access, this approach contributes to the improved performance observed in our benchmarks. Furthermore, the native implementation of core geometric transformations in Julia enables MedImages.jl to maintain the automatic differentiation graph across processing steps, which is critical for the next generation of physics-informed models in molecular imaging and quantitative theranostics.

\section{Conclusions}
The development of MedImages.jl addresses core computational and metadata challenges in modern medical imaging, helping to bridge the AI chasm in nuclear medicine. By leveraging Julia's compilation architecture and type system, the library provides a unified framework where physical accuracy, computational performance, and end-to-end differentiability are achieved simultaneously. The native implementation of spatial transformations allows for transparent integration with SciML pipelines, offering a robust foundation for precision care and ``theranostic digital twins.'' 

Crucially, MedImages.jl resolves the historical tension between dosimetric inaccuracy and clinical intractability: while a gold-standard Monte Carlo simulation requires \textbf{48 hours} of computation to produce a single dose map <Kim_2024_CoRX_NET.pdf>, the proposed framework delivers MC-equivalent accuracy at near real-time speeds. By bridging the divide between the physics of image acquisition and the logic of AI inference, MedImages.jl empowers researchers to unlock the full potential of medical imaging as a precise, quantitative science.

\section*{Availability and Future Directions}
MedImages.jl is open-source and available at \url{https://github.com/JuliaHealth/MedImages.jl}. Future work will focus on integrating with advanced reconstruction frameworks and expanding support for differentiable registration algorithms.

\section{Appendix}

Architectural Foundations: The Coordinate System and Metadata
To accurately integrate functional data with anatomical data across different spatial resolutions, software must rigorously maintain the mapping between discrete voxel indices and physical patient space. This mapping from voxel coordinates $\mathbf{v} = [i, j, k, 1]^\top$ to continuous world coordinates $\mathbf{p} = [x, y, z, 1]^\top$ relies on a $4 \times 4$ homogeneous affine transformation matrix $M$:
\begin{equation}
\mathbf{p} = M \mathbf{v}, \quad M = 
\begin{bmatrix} 
R S & T \\
\mathbf{0} & 1
\end{bmatrix}
\end{equation}
where $R$ is the $3 \times 3$ rotation matrix, $S$ is a diagonal scaling matrix (voxel spacing), and $T$ is the translation vector (origin). MedImages.jl ensures that any spatial operation $\mathcal{T}$ applied to the voxel grid is mathematically coupled with an update $M_{new} = M_{old} \times M_{\mathcal{T}}$, preventing spatial desynchronization.

Furthermore, quantitative nuclear medicine relies on precise temporal and clinical metadata. The Standardized Uptake Value (SUV), a critical metric for PET/CT, is calculated as:
\begin{equation}
SUV = \frac{C(t)}{ID / BW}
\end{equation}
where $C(t)$ is the decay-corrected radioactive concentration, $ID$ is the injected dose, and $BW$ is patient body weight. Discarding parameters such as \textit{Acquisition Time} or \textit{Patient Weight} compromises the quantitative utility of the image. MedImages.jl intrinsically protects these fields within its data structures.

\section{Methods}

\subsection*{Software Architecture}
MedImages.jl implements a metadata paradigm inspired by the Brain Imaging Data Structure (BIDS) standard, organizing and validating clinical metadata rigorously. The core \texttt{MedImage} struct enforces metadata integrity through Julia's type system. Unlike array-based deep learning workflows where converting to tensors discards spatial context, MedImages.jl intrinsically binds metadata (origin, spacing, direction, patient ID, acquisition timestamp) to the structural definition of the medical image.

ITK (C++) is used via wrapper primarily for I/O robustness (NIfTI, DICOM parsing), while spatial transformations are implemented natively in Julia. This ensures: (1) \textbf{Differentiability}—operations integrate with Zygote.jl/Enzyme.jl; (2) \textbf{Composability}—seamless integration with Lux.jl and DifferentialEquations.jl.

To support high-throughput analysis of multimodal studies, MedImages.jl introduces \texttt{BatchedMedImage}, which stores 4D voxel data $(x, y, z, \text{batch})$. Crucially, the spatial metadata (origin, spacing, direction) is vectorized. This design allows each 3D slice within the batch to maintain unique spatial properties while residing in a unified 4D tensor, facilitating efficient batched operations where identical geometric transformations can be dispatched to the GPU in a single pass. To minimize I/O bottlenecks for large datasets, the framework supports HDF5-based storage.

\subsubsection*{Code Example: Typical Workflow}
The following snippet demonstrates loading a multimodal pair, rotating them simultaneously using the batched structure, and computing the SUV:

\begin{lstlisting}[basicstyle=\ttfamily\footnotesize]
using MedImages

# Load PET and CT (DICOM metadata is preserved)
pet = load_image("patient_1/pet.dcm")
ct = load_image("patient_1/ct.nii")

# Create a batched tensor to process simultaneously
batch = create_batched_medimage([pet, ct])

# Rotate both volumes by 45 degrees around Z-axis (GPU supported)
rotated_batch = rotate_mi(batch, 45.0)

# Extract SUV statistics for the PET scan using an ROI mask
suv_stats = calculate_suv_statistics(rotated_batch[1], roi_mask)
println("Mean SUV: ", suv_stats.mean_suv)
\end{lstlisting}

\subsection*{Functionalities and I/O Robustness}
MedImages.jl provides a comprehensive set of functionalities optimized for 3D medical image processing:
\begin{itemize}
    \item \textbf{Robust I/O:} Loading of various medical image formats (NIfTI, DICOM) using \texttt{ITKIOWrapper}.
    \item \textbf{Native Julia Transformations:} Functions to modify spatial properties of 3D volumes (\texttt{rotate\_mi}, \texttt{translate\_mi}, \texttt{scale\_mi}, \texttt{resample\_to\_spacing}, \texttt{crop\_mi}, \texttt{pad\_mi}).
    \item \textbf{Resample to Target:} Resamples a moving image to match the spatial geometry of a reference image.
    \item \textbf{GPU Acceleration:} All transformations transparently support GPU execution via KernelAbstractions.jl.
    \item \textbf{HDF5 Support:} Fast loading and saving of 3D images using HDF5, speeding up I/O for intermediate processing and reducing the "loading penalty" common in deep learning pipelines.
    \item \textbf{Autodifferentiation:} All spatial transformations have well-defined gradients via Zygote.jl and Enzyme.jl.
    \item \textbf{Orientation Management:} Robust tools to handle and convert between neuroimaging coordinate systems.
\end{itemize}

\subsection*{Experimental Setup: Evaluating the Four Challenges}
To systematically evaluate the framework against the four core challenges, several experimental pipelines were designed. To ensure transparency in our evaluation methodology, we summarize the specific datasets, sample sizes, and sources utilized across the experimental series in Table~\ref{tab:dataset_mapping}.

\begin{table}[H]
\centering
\caption{Detailed information - Experimental Series 2: the batched multi-modal theranostic processing}
\label{tab:method_info_batched_Lu177}
\begin{tabular}{|p{0.68\textwidth}|p{0.24\textwidth}|}
\hline
\textbf{Methodology} & \textbf{Study-specific information} \\
\hline
Number of patients / datasets used (single case vs.\ multiple cases) & Single patient case (4 modalities) \\
\hline
Source and acquisition details of the imaging data (scanner type, reconstruction parameters) & $[^{177}\text{Lu}]\text{Lu-PSMA}$ SPECT (AC/NAC), anatomical CT, and a voxel-wise absorbed dose map \\
\hline
Exact voxel dimensions and spatial resolution of each modality before transformation & Standardized target grids of $256 \times 256 \times 256$ and $64 \times 64 \times 64$ voxels \\
\hline
Whether modalities were pre-registered or inherently aligned before batching & Resampled to the common SPECT reference grid prior to batching \\
\hline
Definition of the rotation (angle, axis, interpolation method) & $45^\circ$ around the Z-axis, plus a 10-voxel translation along the X-axis \\
\hline
Interpolation schemes used per modality (e.g., linear, nearest-neighbor) & Linear interpolation (\texttt{Linear\_en}) for continuous image data \\
\hline
Handling of differing voxel grids during batched transformation (internal resampling strategy) & \texttt{resample\_to\_image} function applied to conform differing spatial bounds \\
\hline
Quantitative evaluation metrics, if any, beyond visual inspection & L1 Loss, Normalized Cross-Correlation (NCC), and Mean Squared Error (MSE) \\
\hline
Software and hardware environment (CPU/GPU, Julia version) & Julia \texttt{Lux.jl} ecosystem with NVIDIA GPU acceleration via \texttt{CUDA.jl} \\
\hline
Reproducibility details (random seeds, code availability, parameter settings) & Random seed set to 42, learning rate $1\text{e-}3$, available in MedImages.jl \\
To systematically evaluate the framework against the four core challenges, several experimental pipelines were designed.

\subsubsection*{Evaluating Volume and Speed}
For large-scale studies requiring standardized preprocessing of heterogeneous medical volumes, we developed a fused affine preprocessing pipeline that collapses multiple spatial transformations---orientation normalization, spacing resampling, and spatial centering---into a single $4 \times 4$ affine matrix. This matrix is precomputed on the CPU from spatial metadata, then applied via a custom KernelAbstractions.jl kernel with unrolled trilinear interpolation on the GPU.

To evaluate the framework under high-throughput production loads (Challenge 1: Volume), we benchmarked a 100-case preprocessing pipeline on the autoPET dataset \cite{Dexl_2025} against MONAI's \texttt{PersistentDataset} mechanism. A 100-case dataset was selected to represent a typical multimodality workflow bottleneck: physically disparate modalities (PET and CT) must be loaded and strictly aligned to a common spatial reference frame---meaning they share the identical voxel spacing, direction cosines, and origin---so they can be safely concatenated into a single multi-channel tensor for deep learning ingestion. 

To assess raw execution speed (Challenge 2: Speed), we benchmarked isolated transformations across multiple scales. SimpleITK benchmarks were executed strictly on the CPU, as it does not natively support GPU acceleration for these geometric filters. Conversely, MedImages.jl and MONAI benchmarks utilized GPU execution to demonstrate the high throughput enabled by natively fused kernels. All experiments described are open-source and reproducible via the codebase provided in this repository.

\subsubsection*{Evaluating Differentiability (SciML Integrations)}
To validate end-to-end differentiability (Challenge 3), we implemented three paradigms:
\begin{enumerate}
    \item \textbf{Learned Inverse Rotation (Differentiable Geometric CNN):} A neural network learns to predict and invert unknown 3D rotations applied to volumetric data. To avoid confusion with modern attention-based architectures, we clarify that this uses a standard 3D Convolutional Neural Network (CNN) as a spatial regressor. The CNN predicts three Euler angles, which are converted to an affine matrix to generate a differentiable sampling grid for MedImages.jl's \texttt{interpolate\_pure} function. 
    \item \textbf{Organ-Specific Affine Registration:} A 3D CNN predicts per-organ affine parameters. The loss computation is a single GPU kernel fusing affine transformation, a barycenter distance penalty, and trilinear interpolation loss, differentiated via \texttt{Enzyme.autodiff}.
    \item \textbf{Quantitative Dosimetry Refinement via Universal Differential Equations:} A pipeline designed to correct approximate 3D SPECT reconstructions into high-accuracy, patient-specific dose maps. By embedding neural residual networks within a mechanistic differential equation \cite{Berman_2024}, this approach aims to emulate the accuracy of full Monte Carlo (MC) tracking while maintaining the interpretability essential for medical practice.
\end{enumerate}

The clinical imperative for high-accuracy dosimetry has never been more acute as radiation oncology transitions toward personalized interventions like Targeted Radionuclide Therapy (TRT) and Transarterial Radioembolization (TARE). Historically, dosimetry relied on simplified Medical Internal Radiation Dose (MIRD) formalisms that assumed uniform activity distribution within source organs. While standard clinical methods have transitioned to voxel-based S-values and convolution kernels, they often assume a homogeneous medium, leading to significant errors at tissue interfaces or highly heterogeneous regions (e.g., lungs). While Monte Carlo techniques solve this by tracking stochastic particle interactions, their massive computational burden restricts their clinical utility.

To address this, our pipeline models dose deposition as a parameterized dynamical system. Unlike purely data-driven networks that act as "black boxes," UDEs explicitly satisfy the need for interpretability by treating dose refinement over integration steps as an Initial Value Problem (IVP):
\begin{equation}
\frac{dD}{dt} = f_{\text{mechanistic}}(S, C, D) + \text{NN}(S, C, D; \theta)
\end{equation}
Here, $S$ is the measured SPECT activity, $C$ is the CT-derived tissue density, and $D$ is the cumulative dose. The $f_{\text{mechanistic}}$ term provides a strong inductive bias by hard-coding well-understood primary physical laws—specifically, local energy deposition modeled as proportional to activity and local tissue density. This constraint ensures data efficiency and safety, as the model produces physically plausible results even in regions poorly represented in training data. 

The neural residual network, $\text{NN}_{\theta}$, is the "universal" component embedded directly within the differential equation to learn the highly complex, non-local scattering phenomena (e.g., Compton and Rayleigh scattering) and secondary emissions (e.g., Bremsstrahlung) that cause energy to be deposited at a distance from the source activity. By training against high-fidelity MC ground truths, the network learns to approximate the residual difference between the simple mechanistic physics model and the true dose distribution. To ensure strict differentiability across the entire 5D tensor without encountering Zygote.jl array mutability tracking limitations, this IVP is solved using a custom, out-of-place manual 5-step Euler integration loop rather than standard \texttt{SciMLSensitivity} adjoints.

All convolutional neural networks and Universal Differential Equation models in these experiments were optimized using the Adam optimizer with an initial learning rate of $1 \times 10^{-3}$. The architecture for the neural components consists of a 3D CNN with a width of 32 channels and a depth of 3 residual blocks, employing $3 \times 3 \times 3$ convolutional kernels with unit padding. The models were trained over 20 epochs utilizing an 80/20 train-validation split to monitor convergence, with a composite loss function defined as a weighted sum of 70\% Mean Absolute Error (MAE) and 30\% spatial gradient loss to ensure both intensity accuracy and structural preservation.

\section{Results}
To demonstrate that MedImages.jl resolves the core computational and metadata challenges, our results are structured along the axes defined in the Introduction.

\subsection*{Challenge 1: Volume -- Scaling to 100 Cases (HDF5 vs. Caching)}
To evaluate the framework under high-throughput loads, we benchmarked a 100-case preprocessing pipeline on the autoPET dataset. The core computational task was a multi-modal alignment pipeline, resampling highly heterogeneous PET and CT volumes to a standardized $512 \times 512 \times 512$ isotropic grid. 

In clinical practice, "alignment" implies mathematically mapping the voxel grid of a functional image (e.g., PET, which often has lower resolution and varying spatial origins) to the exact coordinate space of an anatomical reference image (e.g., CT). This requires generating a continuous interpolation grid based on the source metadata (origin, spacing, and direction cosine matrix) and resampling the raw voxel intensities. This alignment is the foundational prerequisite for modern multimodality analysis; Convolutional Neural Networks (CNNs) require multi-channel tensors where pixel $(i,j,k)$ in the PET channel corresponds to the exact same physical tissue location as pixel $(i,j,k)$ in the CT channel. Due to the massive scale of 3D modalities, this continuous spatial resampling is often the dominant bottleneck in deep learning pipelines.

For a fair comparison of computational overhead in this high-volume scenario, we evaluated performance for high-frequency training loops (Epochs 2+), assuming data has been initially processed and cached. We compared MedImages.jl utilizing HDF5 against MONAI's \texttt{PersistentDataset} mechanism. It should be noted that this compares different data persistence strategies—HDF5 vs. Pickle/Pt-based caching—as much as the underlying framework performance.

\begin{table}[H]
\centering
\caption{Definitive Performance Comparison (100 Subjects, Epochs 2+)}
\begin{tabular}{|l|c|c|c|}
\hline
\textbf{Phase (per Subject)} & \textbf{MedImages (HDF5)} & \textbf{MONAI (PersistentDataset)} & \textbf{Speedup} \\
\hline
Stage 1: Load from Cache & \textbf{~40 ms} & ~150 ms & 3.7$\times$ \\
Stage 2: Transform (Compute) & \textbf{~50 ms} & ~500 ms & 10.0$\times$ \\
\hline
\textbf{Total Turnaround} & \textbf{~90 ms} & \textbf{~650 ms} & \textbf{7.2$\times$} \\
\hline
\end{tabular}
\label{tab:scaling_comparison}
\end{table}

Overall, MedImages.jl achieved a 7.2$\times$ total turnaround speedup compared to the MONAI PersistentDataset baseline.

\subsection*{Challenge 2: Speed -- GPU Acceleration Benchmarks}
MedImages.jl was benchmarked on CPU and GPU backends (NVIDIA RTX 3090, Intel Core Ultra 9 285K). For $256 \times 256 \times 128$ volumes, MedImages.jl on GPU achieved significant speedups compared to SimpleITK on CPU:
\begin{itemize}
    \item \textbf{Resampling (Nearest Neighbor):} 115$\times$ speedup vs CPU, 17$\times$ vs SimpleITK.
    \item \textbf{Orientation Changes:} 71$\times$ speedup vs CPU, 31$\times$ vs SimpleITK.
    \item \textbf{Fused Affine Transformation:} 135$\times$ speedup vs CPU, 8.1$\times$ vs SimpleITK.
\end{itemize}

Regarding the Fused Affine Pipeline defined in the methods, benchmarking on 100 subjects resampled to a $512^3$ target grid demonstrated per-subject GPU kernel times of under 200\,ms, confirming the pipeline's suitability for biobank-scale processing.

\subsection*{Challenge 3: Differentiability and SciML Validation}
The Julia ecosystem is uniquely positioned for SciML because its automatic differentiation frameworks (Zygote.jl, Enzyme.jl) can differentiate through arbitrary native code, including custom spatial transformations. Nuclear medicine imaging is inherently quantitative, but its key readouts (e.g., SUV, time–activity curves, absorbed dose) are highly sensitive to geometric preprocessing such as resampling, rotation, and registration. In practice, these operations must often be integrated into AI pipelines for motion correction, longitudinal alignment, and learned preprocessing, where transform parameters are optimized jointly with network weights. However, many commonly used imaging backends execute spatial resampling outside the automatic differentiation graph, preventing end-to-end training through geometric steps \cite{Berman_2024}. To demonstrate that MedImages.jl enables fully differentiable volumetric spatial operations suitable for nuclear medicine workflows, we implement several SciML validation tests.

\subsubsection*{Differentiability Verification Tests}
To rigorously validate correctness, we verified that Zygote.jl returns non-trivial, structurally correct gradients across all spatial transforms. For identity-preserving operations (\texttt{translate\_mi}, \texttt{pad\_mi}), the gradient propagates as $\mathbf{1}$ without corruption. For spatially selective operations (\texttt{crop\_mi}), the gradient is $1$ inside the cropped region and $0$ outside. For the fused affine interpolation kernel, we validated CPU/GPU gradient consistency to within $10^{-4}$.

\subsubsection*{Learned Spatial Transformations}
In the Learned Inverse Rotation experiment, end-to-end differentiability of the geometric pipeline was confirmed using synthetic single-case batched configurations. Over 20 training iterations, the CNN successfully optimized the predicted Euler angles, consistently reducing the mean squared reconstruction error (MSE) between the de-rotated output and the original target by over 65\%. This demonstrates that differentiable augmentation and dynamic learned preprocessing---which are often difficult or rare in standardized imaging pipelines \cite{Shi_2023}---can be seamlessly achieved using native Julia operations.

Similarly, the Organ-Specific Affine Registration pipeline converged using the custom fused GPU loss kernel. The ability of Enzyme.jl to compute exact reverse-mode gradients through the highly branched KernelAbstractions.jl code validated that complex spatial regularizations can be trained natively on the GPU without hand-writing CUDA adjoints.

\subsubsection*{Quantitative Dosimetry Refinement Performance}
The UDE-based dosimetry refinement pipeline successfully corrected approximate SPECT reconstructions into high-accuracy equivalents. The Universal Differential Equation model achieved steady convergence, reaching an MSE of $4.14 \times 10^{-4}$. 

Crucially, the interpretability of the UDE approach is profound: it successfully isolates the known mechanistic physics (primary deposition) from the unknown statistical noise (complex scattering). This allows clinical physicists to independently verify the contribution of the mechanistic term, ensuring that primary dose components adhere to established standard models. Simultaneously, the neural residual robustly corrects secondary scattering artifacts rather than hallucinating primary energy distributions. 

This level of precision is paramount in therapeutic nuclear medicine, where ablative doses are delivered to tumor tissue (e.g., Hepatocellular Carcinoma via Yttrium-90 microspheres) while attempting to spare healthy organs. Standard dose kernels fail to account for the heterogeneous density of tissue interfaces and "leakage" across boundaries. By learning these non-local interactions from MC data, the UDE framework generates a more accurate spatial representation of the dose distribution. This enables clinicians to confidently rely on critical volumetric coverage metrics, such as Dose-Volume Histograms (DVHs) to calculate D95 (minimum dose covering 95\% of the target volume) and D70, predicting pathological response rates accurately without the prohibitive computational costs of full Monte Carlo execution.

\subsection*{Challenge 4: Metadata Management and Clinical Fidelity}
We verified the quantitative accuracy of MedImages.jl by comparing its Standardized Uptake Value (SUV) calculations against the gold standard SimpleITK. Due to Julia's high-precision handling of temporal and clinical metadata, MedImages.jl achieves identical results for SUVbw (Body Weight) calculations, matching SimpleITK's output up to $10^{-14}$ in double precision. This ensures that the massive speedups achieved through GPU acceleration do not compromise clinical validity, a fact verifiable through the open-source benchmarks in this repository.

Furthermore, we conducted an end-to-end consistency experiment on the \texttt{autoPET} dataset \cite{Dexl_2025} to verify that metadata preservation and quantitative fidelity are maintained throughout complex spatial transformation pipelines. The evaluation involved a multi-modal sequence consisting of (1) converting PET intensities to SUV, (2) resampling the PET and label masks to the native CT target space, (3) changing the spacing to a 2.0\,mm isotropic grid, (4) applying a 45-degree rotation around the Z-axis, and (5) applying a 10-voxel translation along the X-axis. For the segmentation masks, nearest-neighbor interpolation was strictly applied to prevent boundary blurring, while linear interpolation was utilized for the continuous PET/CT intensities.

Quantitative analysis of the Mean SUV and total Lesion Volume inside the mask demonstrated remarkable consistency before and after these compound transformations. The original Mean SUV was 3.4913 with a volume of 52.55 $\text{cm}^3$, and the final transformed volume exhibited a Mean SUV of 3.4660 and a volume of 51.86 $\text{cm}^3$. The minor deviations in Mean SUV (0.73\%) and Volume (1.32\%) are fundamentally attributable to grid resampling artifacts and interpolation rather than underlying metadata deterioration. Figure \ref{fig:suv_consistency} visually illustrates the multi-modal alignment before and after the full transformation pipeline.

\begin{figure}[H]
    \centering
    \includegraphics[width=\textwidth]{figures/suv_consistency.png}
    \caption{\textbf{Quantitative Consistency Under Compounded Spatial Transformations (Sagittal Views).}
    The four-panel composite shows sagittal slices of the same subject before and after a sequence of spatial transformations (isotropic 2.0\,mm resampling, 45\textdegree{} Z-rotation, X-translation):
    \textbf{Top-left}: CT after transformations;
    \textbf{Top-right}: PET/SUV after transformations;
    \textbf{Bottom-left}: CT in original native space;
    \textbf{Bottom-right}: PET/SUV in original native space.
    MedImages.jl preserves spatial alignment and clinical metadata throughout, with Mean SUV (3.4913 vs 3.4660) and gross lesion volume (52.55 $\text{cm}^3$ vs 51.86 $\text{cm}^3$) varying by $<1.5\%$ across the full pipeline.}
    \label{fig:suv_consistency}
\end{figure}

\subsubsection*{Illustrative Example: Multi-Modal Batched Processing in Theranostics}

Theranostic radioligand therapy workflows integrate longitudinal patient imaging acquired at different time points, on different devices and reconstruction pipelines---e.g., [$^{177}$Lu]Lu-PSMA SPECT with and without attenuation correction (AC/NAC), voxel-wise absorbed dose maps, and CT---where each modality retains its own voxel grid, datatype, and spatial metadata (spacing, origin, orientation).  

To demonstrate the capability of \texttt{BatchedMedImage} in handling complex, multi-modal datasets typical for theranostics workflows, we processed a dataset of four typical examinations, each containing a distinct imaging modality:
\begin{enumerate}
    \item [$^{177}$Lu]Lu-PSMA SPECT with attenuation correction (AC)
    \item [$^{177}$Lu]Lu-PSMA SPECT without attenuation correction (NAC)
    \item [$^{177}$Lu]Lu-PSMA Dosemap (voxel-wise absorbed dose)
    \item Computed Tomography (CT) for anatomical reference
\end{enumerate}

The images were loaded into a single \texttt{BatchedMedImage} object, preserving their individual spatial metadata and voxel types. A batched rotation transformation was then applied simultaneously to all four modalities. Figure \ref{fig:batched_rotation} illustrates the results, showcasing the original axial slices (left) and the rotated output (right). 

The ability to apply identical geometric transformations across functional (SPECT, Dosemap) and anatomical (CT) data in a single operation ensures perfect spatial alignment is maintained during augmentation or registration tasks, which is critical for accurate dosimetry and treatment planning.

\begin{figure}[H]
    \centering
    \begin{minipage}{0.48\textwidth}
        \centering
        \includegraphics[width=\linewidth]{figures/orig.png}
        \caption*{A: Original}
    \end{minipage}
    \hfill
    \begin{minipage}{0.48\textwidth}
        \centering
        \includegraphics[width=\linewidth]{figures/rot.png}
        \caption*{B: Rotated}
    \end{minipage}
    \caption{\textbf{Simultaneous Batched Rotation of Multi-Modal Theranostic Data.} 
    Panel A displays the original axial slices of four distinct modalities stored within a single \texttt{BatchedMedImage} structure: (Top-Left) [$^{177}$Lu]Lu-PSMA SPECT AC, (Top-Right) [$^{177}$Lu]Lu-PSMA SPECT NAC, (Bottom-Left) [$^{177}$Lu]Lu-PSMA Dosemap, and (Bottom-Right) CT. 
    Panel B shows the result after applying a single \texttt{rotate\_mi} operation to the batch. The spatial integrity and alignment between the functional distribution (SPECT/Dose) and anatomical structure (CT) are preserved during the transformation.}
    \label{fig:batched_rotation}
\end{figure}


\section{Discussion}
As biomedical imaging increasingly embraces complex multimodal, multidimensional, and longitudinal studies, alongside scientific machine learning (SciML), the interoperability and performance of data management frameworks become critical. Our results indicate that MedImages.jl’s metadata handling paradigm, combined with seamless integration into the Julia scientific ecosystem, provides a highly performant and differentiable framework that systematically addresses the three core challenges outlined in the Introduction.  

\subsection*{Addressing Challenges 1 and 2 (Volume and Speed): The Advantage of Native Fused Kernels}
Integration of metadata within the \texttt{MedImage} structure directly translates into performance gains observed in our benchmarks. The 7.2$\times$ turnaround advantage observed in high-frequency training loops is primarily attributed to MedImages.jl's integration with HDF5 -- which bypasses the serialization costs inherent in many Python caching mechanisms -- and its single-pass Fused Affine kernel. By mathematically collapsing the entire spatial transformation stack (normalization, resampling, centering, etc.) into a single CUDA execution, MedImages.jl eliminates the memory traffic and synchronization overhead associated with sequential operation queues, providing a highly efficient solution for biobank-scale deployments.

\subsection*{Addressing Challenge 2 (Metadata Management): Comparing Metadata Architectures}
Modern Python frameworks like MONAI \cite{CardosoLi2022} have introduced sophisticated mechanisms to couple metadata with image arrays. MONAI utilizes a MetaTensor structure, appending a metadata dictionary to the PyTorch Tensor. A similar design pattern appears in TorchIO \cite{perezgarcia2021torchio}, which prioritizes convenient multimodal coordination but still builds its core processing around external imaging backends. TorchIO approaches the problem through a hierarchical Subject class system. Although effective for multi-modal synchronization, TorchIO is primarily a manager of third-party backends  like SimpleITK or NiBabel cite\. This the dependency chain maintains a layer of indirection that can lead to memory management issues in large-scale training \cite{perezgarcia2021torchio}.

In contrast, MedImages.jl leverages Julia's type system to bind metadata directly within the \texttt{MedImage} struct. As summarized in Table \ref{tab:metadata_comparison}, by removing the runtime overhead associated with dictionary-based metadata access, this approach contributes to the improved performance observed in our benchmarks. Furthermore, the native implementation of core geometric transformations in Julia enables MedImages.jl to maintain the automatic differentiation graph across processing steps. This is particularly relevant for imaging workflows that incorporate scientific machine learning, where gradient information may be lost when operations are passed to non-differentiable compiled backends cite\. 

\begin{table}[H]
\centering
\caption{Comparative Metadata Architectures in Modern Frameworks}
\begin{tabular}{|p{2.5cm}|p{3.5cm}|p{3.5cm}|p{3.5cm}|}
\hline
\textbf{Framework} & \textbf{Metadata Storage Mechanism} & \textbf{Preservation Strategy} & \textbf{Potential Interoperability Constraints} \\
\hline
SimpleITK & Internal C++ Struct & Encapsulated in \texttt{itk::Image} & User-dependent retention (often discarded upon conversion to NumPy). \\
\hline
MONAI & MetaTensor (Dict-based) & Metadata dictionary moves with tensor & Dictionary lookup overhead; potential orientation flip during ITK interoperability. \\
\hline
TorchIO & Subject Class Attributes & Multi-modal synchronization & Potential memory buildup during extensive transformation histories. \\
\hline
MedImages.jl & Typed Julia Struct & Explicit binding within struct & N/A (Native Julia stack). \\
\hline
\end{tabular}
\label{tab:metadata_comparison}
\end{table}

\subsection*{The Primacy of Clinical and Temporal Metadata}
The lack of standardization in medical imaging preprocessing pipelines can pose a reproducibility challenge \cite{Esteban_2018}. In such settings, metadata drift, defined as the decoupling of spatial or clinical metadata from voxel data, may introduce clinically relevant errors. For example, it can lead to displacement of treatment margins in radiotherapy or compromise quantitative interpretation in PET/CT when decay-related timing information is lost \textbackslash{}cite. Quantitative nuclear medicine is particularly dependent on the integrity of clinical and temporal metadata, as standardized uptake value (SUV) calculations explicitly require accurate injection and acquisition times. In theranostic workflows, this requirement must be maintained while applying identical geometric transformations across multiple co-registered modalities, such as [$^{177}$Lu]Lu-PSMA SPECT AC/NAC, absorbed-dose maps, and CT.

MedImages.jl addresses this by design: the \texttt{BatchedMedImage} abstraction allows single-call, batched geometric transforms to be applied consistently across modalities without manual bookkeeping. Our cross-platform validation confirms that the SUV conversion factors extracted by MedImages.jl achieve mathematical identity with established SimpleITK benchmarks to a precision of $10^{-14}$. This numerical equivalence ensures that researchers transitioning to Julia can maintain absolute continuity with legacy clinical findings while gaining the performance required for massive biobank-scale analysis.

\subsection*{Addressing Challenge 3 (Differentiability): Navigating the Framework Ecosystem}
We do not advocate replacing established tools wholesale. Rather, we recommend choosing the framework that best matches the task at hand, especially when differentiability and complex scientific computing are core requirements.

\paragraph{When established Python frameworks remain the pragmatic choice.}
For teams with substantial legacy Python codebases, switching frameworks carries non-trivial migration costs. MONAI provides a mature library of ready-made components---pretrained segmentation networks, standardized data-loading pipelines, and extensive augmentation transforms---that can be deployed with minimal custom code. Similarly, SimpleITK offers a comprehensive catalogue of classical filters whose correctness has been validated across decades. When the research question can be answered with an existing pipeline and the primary requirement is rapid prototyping within a specific tensor framework (like PyTorch), these tools remain highly efficient.

\paragraph{When MedImages.jl offers a clear advantage for SciML.}
The Julia-based approach becomes compelling when research transcends standard convolutional network training and enters the realm of physics-informed models or custom numerical workflows \cite{RoeschGreener2023}:
\begin{itemize}
    \item \textbf{Unifying the Fractured SciML Landscape.} In Python, moving a differentiable model from PyTorch to a JAX-based differential equation solver often requires significant code rewriting or bridging overhead. Julia's multiple dispatch solves the expression problem, meaning essentially all packages are interoperable natively \cite{PalBhattacharya2024, EschleGl2023}. MedImages.jl serves as the critical entry point to make this uniquely unified SciML ecosystem accessible for medical imaging.
    \item \textbf{Novel algorithm development.} When researchers need to implement a new spatial transformation, loss function, or physics-informed model from scratch, Julia's single-language stack eliminates the need to write C++/CUDA extensions. The algorithm can be prototyped, debugged, and GPU-accelerated within the same language, with full automatic differentiation support.
    \item \textbf{Cross-domain scientific computing.} Julia's ecosystem uniquely bridges medical imaging with adjacent scientific domains---DifferentialEquations.jl for pharmacokinetic modelling, KomaMRI.jl \cite{CastilloPassiCoronado2023} for MRI pulse-sequence simulation, Flux.jl/Lux.jl \cite{pal2023lux} for deep learning---all within a single runtime. This composability enables workflows that would otherwise require fragile inter-process communication between Python, MATLAB, and C++ components.
    \item \textbf{Differentiable imaging physics.} Tasks such as learned registration, differentiable augmentation, or physics-informed neural networks require gradients to flow through spatial transformations. MedImages.jl's native Julia kernels integrate directly with Zygote.jl and Enzyme.jl, whereas MONAI and TorchIO must fall back on non-differentiable C++ backends for many geometric operations.
\end{itemize}


MedImages.jl offers a path forward---a future where imaging physics is integrated directly into the deep learning pipeline, where hardware acceleration is transparent and ubiquitous, and where the metadata that defines our physical world is preserved with the same rigor as the pixel values themselves. By bridging the divide between the physics of acquisition and the logic of inference, and by addressing the entwined challenges of \textbf{Volume} and \textbf{Speed} with a 7.2$\times$ faster turnaround than traditional caching frameworks \cite{Shi_2023}, MedImages.jl empowers researchers and clinicians to unlock the full potential of medical imaging as a precise, quantitative science. This performance is especially critical for large-scale multi-center studies and meta-analyses, where the ability to process thousands of subjects with absolute metadata integrity enables research that was previously computationally infeasible.


\subsection*{Limitations}
While MedImages.jl offers significant advantages in execution speed and differentiability, adopting the Julia ecosystem presents certain trade-offs. The most notable is the "time-to-first-plot" or compilation latency. Because Julia uses Just-In-Time compilation, the first execution of a complex spatial transformation incurs a compilation penalty, making initial startup slower than equivalent interpreted Python code. Furthermore, running full test suites or deploying models with heavy dependencies (e.g., CUDA.jl, Lux.jl, Zygote.jl) requires significant development memory, occasionally causing out-of-memory (OOM) issues on constrained hardware. Finally, while the Julia machine learning ecosystem is growing rapidly, it does not yet match the sheer volume of pre-trained models and plug-and-play architectural components available in the PyTorch/MONAI ecosystem.

\section{Conclusions}
The development of MedImages.jl addresses core computational and metadata challenges in modern medical imaging. By leveraging Julia's compilation architecture and type system, the library provides a unified framework where physical accuracy, computational performance, and end-to-end differentiability are achieved simultaneously. The native implementation of spatial transformations allows for transparent integration with SciML pipelines, offering a robust foundation for the next generation of physics-informed models in molecular imaging and quantitative theranostics.
In practice, we envision a complementary workflow: established MONAI/SimpleITK pipelines for routine clinical deployment, and MedImages.jl for research frontiers where performance, differentiability, or algorithmic novelty are the primary requirements.

\section*{Availability and Future Directions}
MedImages.jl is open-source and available at \url{https://github.com/JuliaHealth/MedImages.jl}. The results and architecture described in this manuscript correspond to MedImages.jl version v2.0.1. All experiments described in this paper, including performance benchmarks and differentiability validations, are reproducible via the scripts provided in the repository's \texttt{experiment} directory.

Future work will focus on:
\begin{enumerate}
    \item Integrating with advanced reconstruction frameworks (KomaMRI.jl \cite{CastilloPassiCoronado2023}, GIRFReco.jl \cite{JaffrayWu2024}) for end-to-end physics-informed learning.
    \item Implementing differentiable registration algorithms leveraging Zygote.jl.
    \item Expanding support for standardized formats like DICOM-SEG via Highdicom \cite{BridgeGorman2022} integration.
    \item Developing architectural foundations for clinical-grade AI agents (e.g., VoxelPrompt) capable of multi-modal "Medical Superintelligence" benchmarks where absolute metadata integrity is a prerequisite for safe automated diagnostics.
\end{enumerate}

\section*{Declarations}
\textbf{Ethics approval and consent to participate:} This research relies solely on retrospective analysis of fully anonymized, publicly available datasets. No new human subjects data was collected. Therefore, formal ethics committee approval and informed consent were not required. \\
\textbf{Declaration of competing interests:} The authors declare that they have no known competing financial interests or personal relationships that could have appeared to influence the work reported in this paper. \\
\textbf{Funding sources:} This research did not receive any specific grant from funding agencies in the public, commercial, or not-for-profit sectors. \\
\textbf{Declaration of generative AI and AI-assisted technologies:} During the preparation of this work, the author(s) used generative AI and AI-assisted technologies for improving style and assisting in code development. After using these tools, the authors reviewed and edited the content as needed and take full responsibility for the content of the published article.

\bibliographystyle{elsarticle-num}
\bibliography{bibl}

\end{document}

\bibliography{bibl}

\end{document}
ontent of the published article.

\bibliographystyle{elsarticle-num}
\bibliography{bibl}

\end{document}
